% Appendix-only theoretical note for an ICLR-style paper.
% Intended use: place after \appendix and include with \input{...}.
%
% Preamble requirements, if not already present:
%   \usepackage{amsmath,amssymb,mathtools,amsthm}
%   \newtheorem{assumption}{Assumption}[section]
%   \newtheorem{lemma}{Lemma}[section]
%   \newtheorem{proposition}{Proposition}[section]
%   \newtheorem{theorem}{Theorem}[section]
%   \newtheorem{corollary}{Corollary}[section]
%   \newtheorem{remark}{Remark}[section]
%
% Optional macros. Remove these if already defined in the main paper.
\providecommand{\rms}{\operatorname{rms}}
\providecommand{\diag}{\operatorname{diag}}
\providecommand{\MP}{\operatorname{MP}}
\providecommand{\AW}{\operatorname{AdamW}}
\providecommand{\E}{\mathbb{E}}
\providecommand{\R}{\mathbb{R}}
\providecommand{\cG}{\mathcal{G}}
\providecommand{\cQ}{\mathcal{Q}}
\providecommand{\cP}{\mathcal{P}}
\providecommand{\norm}[1]{\left\lVert #1 \right\rVert}
\providecommand{\ip}[2]{\left\langle #1,#2 \right\rangle}
\providecommand{\lammax}{\lambda_{\max}}
\providecommand{\lammin}{\lambda_{\min}}

\section{A Conditional Optimization Separation for MatrixPolicy}
\label{app:complete-opt-proof}
\label{app:gauge-conditioning}

This appendix gives a conditional optimization guarantee for MatrixPolicy in the
local regime suggested by the RLB parameterization.  The assumptions are stated
explicitly: the rational block is in a stable local regime, the loss is well
approximated by a strongly convex quotient objective, and optimizer states have
entered a frozen-preconditioner phase.  In that regime, MatrixPolicy has a
strictly better contraction factor than AdamW whenever its canonicalized quotient
preconditioner has a smaller condition number.

The proof has three steps.  First, the RLB block induces a positive scale gauge, so many
matrix pairs represent the same function.  Second, MatrixPolicy's rebalance step selects a
bounded representative of that gauge orbit, whereas AdamW operates in raw coordinates and
therefore induces a gauge-dependent quotient preconditioner.  Third, for a strongly convex
quadratic quotient objective, the optimizer with the smaller quotient condition number has
the smaller worst-case linear convergence factor.  This gives a complete optimization
separation under the stated assumptions.

\subsection{Gauge quotient model}
\label{app:complete-opt-proof:model}

For one RLB layer, suppress the layer index and write
\begin{align}
    z_g &= A_g x, \\
    r_g &= \sqrt{m^{-1}\norm{z_g}_2^2+\epsilon_{\rm rms}}, \\
    u_g &= z_g/r_g, \\
    h_g &= r_g R_g(u_g), \\
    y &= B\operatorname{concat}_{g=1}^G h_g .
\end{align}
For $a_g>0$, define the positive scale action
\begin{equation}
    T_a(A_g,B_g)=(a_gA_g,a_g^{-1}B_g),
    \qquad g=1,\ldots,G .
    \label{eq:complete-positive-action}
\end{equation}
Equivalently, with $s_g=\log a_g$,
\begin{equation}
    T_s(A_g,B_g)=(e^{s_g}A_g,e^{-s_g}B_g).
    \label{eq:complete-log-action}
\end{equation}
When $\epsilon_{\rm rms}=0$, the RLB block is homogeneous and
$T_s(A,B)$ represents the same function as $(A,B)$.  With a nonzero RMS floor, this
symmetry is approximate on nondegenerate activation regions; the perturbation is handled
in Proposition~\ref{prop:complete-perturbation} below.

Let $\cP$ denote the local parameter space of the RLB matrices and let
$\cG=(\R,+)^G$ denote the log-gauge group.  We write the quotient space as
$\cQ=\cP/\cG$.  A point of $\cQ$ represents a function-space degree of freedom, whereas
the gauge coordinate only selects a matrix representative of the same function.

\begin{assumption}[Local quotient chart]
\label{ass:quotient-chart}
There is a neighborhood of the training trajectory in which parameters can be written as
\begin{equation}
    \theta=\theta(q,s),
    \qquad q\in\R^d,
    \qquad s\in\R^G,
\end{equation}
where $q$ is a local coordinate on the quotient $\cQ$ and $s$ is the log-gauge coordinate.
The objective depends on $\theta$ only through $q$:
\begin{equation}
    L(\theta(q,s))=F(q).
    \label{eq:loss-quotient}
\end{equation}
For the exact theorem we take $\epsilon_{\rm rms}=0$.  For the perturbed theorem we assume
that $\epsilon_{\rm rms}$ is small enough that the quotient description holds up to an
$O(\epsilon_{\rm rms})$ error on the local region.
\end{assumption}

\begin{assumption}[Quadratic local objective]
\label{ass:quadratic-objective}
In the quotient chart, the loss is a strongly convex quadratic
\begin{equation}
    F(q)=F_\star+\frac{1}{2}(q-q_\star)^\top H(q-q_\star),
    \label{eq:quadratic-F}
\end{equation}
where $H\succ 0$ and
\begin{equation}
    0<\mu I\preceq H\preceq L I<\infty .
\end{equation}
\end{assumption}

Assumption~\ref{ass:quadratic-objective} is the standard local model behind
condition-number analyses of first-order methods and gives a complete closed-form
contraction proof.  A smooth strongly convex objective
can be reduced to this setting locally by replacing $H$ with the Hessian at $q_\star$ and
controlling the Taylor remainder; the finite-perturbation version is stated in
Proposition~\ref{prop:complete-perturbation}.

\subsection{MatrixPolicy canonicalization}
\label{app:complete-opt-proof:canonicalization}

MatrixPolicy uses the gauge freedom to select a controlled representative of the same
quotient point.  For group $g$, define the log scale ratio
\begin{equation}
    \gamma_g=\log \rms(A_g)-\log \rms(B_g).
    \label{eq:complete-gamma}
\end{equation}
Let $\tau_g$ be the target log ratio determined by the optimizer's gain and pressure
statistics.  Ignoring clipping, the rebalance step applies
\begin{equation}
    \ell_g=\frac{1}{2}(\tau_g-\gamma_g),
    \qquad
    A_g\leftarrow e^{\ell_g}A_g,
    \qquad
    B_g\leftarrow e^{-\ell_g}B_g.
    \label{eq:complete-rebalance}
\end{equation}

\begin{lemma}[Exact targeting of the representative]
\label{lem:complete-targeting}
If the clipping in \eqref{eq:complete-rebalance} is inactive, then the rebalance maps
$\gamma_g$ exactly to $\tau_g$.
\end{lemma}

\begin{proof}
After the rebalance,
\begin{align}
    \gamma_g^+
    &=\log \rms(e^{\ell_g}A_g)-\log \rms(e^{-\ell_g}B_g) \\
    &=\gamma_g+2\ell_g
      =\gamma_g+\tau_g-\gamma_g
      =\tau_g .
\end{align}
\end{proof}

Thus, in the homogeneous block, the rebalance changes only the gauge coordinate and leaves
the quotient coordinate $q$ unchanged.  In the exact theorem below we assume the target is
inside the clipping interval.  If clipping, the RMS floor, or finite precision introduces a
small deviation from exact canonicalization, that deviation is included in the perturbation
term of Proposition~\ref{prop:complete-perturbation}.

\subsection{Optimizer-induced quotient dynamics}
\label{app:complete-opt-proof:quotient-dynamics}

A raw matrix optimizer induces a motion on the quotient coordinate $q$.  In a local
frozen-state regime, both MatrixPolicy and AdamW can be written as preconditioned gradient
steps on $F$:
\begin{equation}
    q_{t+1}=q_t-\eta P\nabla F(q_t),
    \label{eq:preconditioned-q-update}
\end{equation}
where $P\succ0$ is the quotient preconditioner induced by the optimizer and by the chosen
gauge representative.  The next assumption makes this statement exact.

\begin{assumption}[Frozen-preconditioner optimizer model]
\label{ass:frozen-preconditioner}
During the local phase analyzed here, the two optimizers induce fixed positive definite
quotient preconditioners $P_{\MP}$ and $P_{\AW}$ satisfying
\begin{align}
    q_{t+1}^{\MP}
    &=q_t^{\MP}-\eta_{\MP}P_{\MP}\nabla F(q_t^{\MP}),
    \label{eq:mp-ideal-dynamics}\\
    q_{t+1}^{\AW}
    &=q_t^{\AW}-\eta_{\AW}P_{\AW}\nabla F(q_t^{\AW}).
    \label{eq:aw-ideal-dynamics}
\end{align}
Here $P_{\MP}$ is the preconditioner after MatrixPolicy has applied its gauge rebalance,
and $P_{\AW}$ is the quotient preconditioner induced by AdamW's raw-coordinate diagonal
adaptive state at the same represented function.  Weight decay is set to zero in the exact
theorem, or else is treated as a perturbation in
Proposition~\ref{prop:complete-perturbation}.
\end{assumption}

This idealized frozen-state model isolates the mechanism: AdamW is analyzed as
the diagonal adaptive method seen by the quotient after its moment estimates have
stabilized, while MatrixPolicy is analyzed after canonicalizing the RLB gauge.

For a preconditioner $P\succ0$, define the quotient Hessian seen by the optimizer as
\begin{equation}
    K(P)=P^{1/2}HP^{1/2},
    \label{eq:quotient-hessian}
\end{equation}
and define its condition number
\begin{equation}
    \kappa(P;H)=\frac{\lammax(K(P))}{\lammin(K(P))}.
    \label{eq:precond-condition-number}
\end{equation}
The best constant step size for the quadratic under preconditioner $P$ is
\begin{equation}
    \eta^\star(P)=\frac{2}{\lammax(K(P))+\lammin(K(P))},
    \label{eq:best-stepsize}
\end{equation}
with contraction factor
\begin{equation}
    \rho(P;H)
    =\frac{\kappa(P;H)-1}{\kappa(P;H)+1}.
    \label{eq:rho-def}
\end{equation}

\begin{assumption}[Strict quotient-conditioning advantage]
\label{ass:conditioning-advantage}
MatrixPolicy's canonicalized quotient preconditioner has a strictly smaller condition
number than AdamW's raw-coordinate quotient preconditioner:
\begin{equation}
    \kappa_{\MP}:=\kappa(P_{\MP};H)
    <
    \kappa(P_{\AW};H)=:\kappa_{\AW} .
    \label{eq:kappa-gap}
\end{equation}
Equivalently, for some $\Delta_\kappa>0$,
\begin{equation}
    \kappa_{\AW}\geq \kappa_{\MP}+\Delta_\kappa .
\end{equation}
\end{assumption}

Assumption~\ref{ass:conditioning-advantage} is the mechanism condition.  In the
local regime being modeled, MatrixPolicy's role-aware matrix step and gauge
rebalance produce a better conditioned quotient update than the AdamW update
induced by the raw matrix coordinates.  This is the formal version of the
mechanism suggested by the empirical comparison.

\subsection{Why the condition-number gap is a gauge effect}
\label{app:complete-opt-proof:gauge-effect}

The condition-number assumption is natural for RLB because the raw Euclidean coordinates
contain a nonfunctional scale direction.  We record the elementary calculation.

\begin{lemma}[Gauge directions are null function directions]
\label{lem:gauge-null-direction}
Assume $\epsilon_{\rm rms}=0$.  For any differentiable loss depending on the RLB matrices
only through the represented function,
\begin{equation}
    \ip{\nabla_{A_g}L}{A_g}-\ip{\nabla_{B_g}L}{B_g}=0 .
    \label{eq:gauge-null-gradient}
\end{equation}
\end{lemma}

\begin{proof}
Since $L(T_s(A_g,B_g))=L(A_g,B_g)$ for all $s_g$, differentiate at $s_g=0$:
\begin{align}
    0
    &=\frac{d}{ds_g}L(e^{s_g}A_g,e^{-s_g}B_g)\bigg|_{s_g=0} \\
    &=\ip{\nabla_{A_g}L}{A_g}-\ip{\nabla_{B_g}L}{B_g}.
\end{align}
\end{proof}

Thus $(A_g,-B_g)$ changes the representative but not the function.  AdamW leaves
this representative dependence in its raw-coordinate diagonal moments, denominator
terms, and decoupled weight decay.  Consequently, two gauge-equivalent parameterizations can
induce different $P_{\AW}$ on the quotient.

The following toy calculation shows the scaling explicitly.  Consider the scalar quotient
coordinate $w=ab$, with gauge $a=e^s\sqrt{w}$ and $b=e^{-s}\sqrt{w}$.  For
$\ell(w)=\frac{1}{2}(w-y)^2$, raw Euclidean gradient descent in $(a,b)$ gives, to first
order in the step size,
\begin{equation}
    w^+
    =w-\eta(a^2+b^2)(w-y)+O(\eta^2).
\end{equation}
Since
\begin{equation}
    a^2+b^2=w(e^{2s}+e^{-2s})=2w\cosh(2s),
\end{equation}
the quotient step size depends on the arbitrary gauge $s$.  MatrixPolicy's
rebalance sets $s$ to a controlled target, while AdamW operates in the raw
representative.  In higher dimensions, the same
calculation appears as a gauge-dependent quotient metric, which is summarized by
$\kappa(P_{\AW};H)$.

\subsection{Contraction theorem}
\label{app:complete-opt-proof:theorem}

We now prove the optimization separation.  The proof is standard for preconditioned
gradient descent on a quadratic, but the notation makes explicit that the relevant
condition number is the quotient condition number, not the raw parameter condition number.

\begin{lemma}[Preconditioned quadratic contraction]
\label{lem:preconditioned-contraction}
Let $F$ satisfy Assumption~\ref{ass:quadratic-objective}.  For any $P\succ0$, consider
\begin{equation}
    q_{t+1}=q_t-\eta^\star(P)P\nabla F(q_t),
\end{equation}
where $\eta^\star(P)$ is defined in \eqref{eq:best-stepsize}.  Then
\begin{equation}
    F(q_T)-F_\star
    \leq
    \rho(P;H)^{2T}\bigl(F(q_0)-F_\star\bigr),
    \label{eq:upper-contraction}
\end{equation}
where $\rho(P;H)$ is defined in \eqref{eq:rho-def}.  Moreover, the bound is sharp: for
each $T$, there exists an initialization $q_0$ for which equality holds.
\end{lemma}

\begin{proof}
Let $e_t=q_t-q_\star$ and define $\widetilde e_t=P^{-1/2}e_t$.  Since
$\nabla F(q_t)=He_t$, the update becomes
\begin{equation}
    \widetilde e_{t+1}
    =\left(I-\eta^\star(P)K(P)\right)\widetilde e_t,
\end{equation}
where $K(P)=P^{1/2}HP^{1/2}$.  Also,
\begin{equation}
    2(F(q_t)-F_\star)=e_t^THe_t=\widetilde e_t^TK(P)\widetilde e_t.
\end{equation}
Diagonalize $K(P)=U\Lambda U^T$ and write $v_t=U^T\widetilde e_t$.  Then
\begin{equation}
    v_{t+1,i}=\left(1-\eta^\star(P)\lambda_i\right)v_{t,i}.
\end{equation}
The largest absolute multiplier is
\begin{equation}
    \max_i\left|1-\eta^\star(P)\lambda_i\right|
    =
    \frac{\lammax(K(P))-\lammin(K(P))}
    {\lammax(K(P))+\lammin(K(P))}
    =\rho(P;H).
\end{equation}
Therefore
\begin{align}
    2(F(q_{t+1})-F_\star)
    &=\sum_i \lambda_i v_{t+1,i}^2 \\
    &\leq \rho(P;H)^2\sum_i \lambda_i v_{t,i}^2 \\
    &=2\rho(P;H)^2(F(q_t)-F_\star).
\end{align}
Iterating gives \eqref{eq:upper-contraction}.  Equality is obtained by choosing
$\widetilde e_0$ to be an eigenvector of $K(P)$ associated with either
$\lammin(K(P))$ or $\lammax(K(P))$.
\end{proof}

\begin{theorem}[Conditional optimization separation]
\label{thm:conditional-optimization-separation}
Suppose Assumptions~\ref{ass:quotient-chart}--\ref{ass:conditioning-advantage} hold, and
choose the locally optimal constant step sizes
\begin{equation}
    \eta_{\MP}=\eta^\star(P_{\MP}),
    \qquad
    \eta_{\AW}=\eta^\star(P_{\AW}).
\end{equation}
Let
\begin{equation}
    \rho_{\MP}=\frac{\kappa_{\MP}-1}{\kappa_{\MP}+1},
    \qquad
    \rho_{\AW}=\frac{\kappa_{\AW}-1}{\kappa_{\AW}+1}.
\end{equation}
Then
\begin{equation}
    \rho_{\MP}<\rho_{\AW}<1 .
    \label{eq:rho-strict-gap}
\end{equation}
Consequently, MatrixPolicy has a strictly smaller worst-case linear convergence factor on
the quotient objective:
\begin{equation}
    \sup_{q_0\neq q_\star}
    \frac{F(q_T^{\MP})-F_\star}{F(q_0)-F_\star}
    \leq
    \rho_{\MP}^{2T}
    <
    \rho_{\AW}^{2T}
    =
    \sup_{q_0\neq q_\star}
    \frac{F(q_T^{\AW})-F_\star}{F(q_0)-F_\star}
    \label{eq:main-separation}
\end{equation}
for every integer $T\geq1$.
\end{theorem}

\begin{proof}
The map $\kappa\mapsto(\kappa-1)/(\kappa+1)$ is strictly increasing for
$\kappa>1$.  Assumption~\ref{ass:conditioning-advantage} therefore implies
$\rho_{\MP}<\rho_{\AW}$.  Since both condition numbers are finite and at least one,
$0\leq\rho_{\MP}<\rho_{\AW}<1$ unless both methods are perfectly conditioned, in which
case the strict condition-number gap cannot hold.  Applying
Lemma~\ref{lem:preconditioned-contraction} to $P_{\MP}$ and $P_{\AW}$ gives the two
worst-case contraction factors.  The strict inequality in \eqref{eq:main-separation}
follows from $\rho_{\MP}<\rho_{\AW}$ and $T\geq1$.
\end{proof}

The theorem compares the best constant learning rate for each local preconditioner, so the
separation is not caused by giving MatrixPolicy a better scalar step size.  It comes only
from the quotient condition number.  In this sense, the result is an optimizer-level
statement: after quotienting the RLB gauge, the method with the better induced quotient
metric converges faster on the local quadratic model.

\subsection{Finite-time dominance for a fixed initialization}
\label{app:complete-opt-proof:fixed-init}

Theorem~\ref{thm:conditional-optimization-separation} is a worst-case statement.  A fixed
initialization version follows by assuming that the AdamW trajectory has a component in an
extremal eigendirection of its quotient Hessian.

\begin{assumption}[Nontrivial AdamW hard-mode component]
\label{ass:hard-mode}
Let $K_{\AW}=K(P_{\AW})$.  In the eigenbasis of $K_{\AW}$, the initial transformed error
$P_{\AW}^{-1/2}(q_0-q_\star)$ has at least an $\alpha$ fraction of its $K_{\AW}$-energy in
eigendirections whose one-step multiplier under AdamW has absolute value at least
$\rho_{\AW}-\xi$, where $\alpha\in(0,1]$ and $0\leq\xi<\rho_{\AW}-\rho_{\MP}$.
\end{assumption}

\begin{corollary}[Fixed-initialization separation]
\label{cor:fixed-init-separation}
Under the assumptions of Theorem~\ref{thm:conditional-optimization-separation} and
Assumption~\ref{ass:hard-mode}, AdamW satisfies
\begin{equation}
    F(q_T^{\AW})-F_\star
    \geq
    \alpha(\rho_{\AW}-\xi)^{2T}\bigl(F(q_0)-F_\star\bigr).
    \label{eq:adam-lower-fixed}
\end{equation}
MatrixPolicy satisfies
\begin{equation}
    F(q_T^{\MP})-F_\star
    \leq
    \rho_{\MP}^{2T}\bigl(F(q_0)-F_\star\bigr).
    \label{eq:mp-upper-fixed}
\end{equation}
Thus, whenever
\begin{equation}
    \rho_{\MP}^{2T}<\alpha(\rho_{\AW}-\xi)^{2T},
    \label{eq:fixed-separation-condition}
\end{equation}
MatrixPolicy has strictly lower quotient loss after $T$ steps.
\end{corollary}

\begin{proof}
Equation~\eqref{eq:mp-upper-fixed} is Lemma~\ref{lem:preconditioned-contraction}.  For
AdamW, decompose the transformed error in the eigenbasis of $K_{\AW}$.  The part of the
initial energy in the hard eigendirections is at least an $\alpha$ fraction by
Assumption~\ref{ass:hard-mode}, and every such component is multiplied by at least
$\rho_{\AW}-\xi$ in absolute value each step.  This gives
\eqref{eq:adam-lower-fixed}.  Combining \eqref{eq:adam-lower-fixed} and
\eqref{eq:mp-upper-fixed} gives the strict separation condition
\eqref{eq:fixed-separation-condition}.
\end{proof}

The additional assumption is mild in high-dimensional quadratic models: unless the initial
error is specially aligned away from AdamW's hard directions, some energy is present in
those directions.  The corollary is nevertheless stated explicitly because fixed-run
comparisons always depend on initialization.

\subsection{Perturbation stability}
\label{app:complete-opt-proof:perturbation}

The exact theorem used a homogeneous radius, frozen optimizer states, zero weight decay,
and a quadratic loss.  The next proposition states the usual finite-time stability fact:
if the implemented dynamics are uniformly close to the ideal dynamics, then a strict
finite-time gap persists.  This is the appropriate $\delta$--$\epsilon$ statement for the
proof.

\begin{assumption}[Uniform local perturbation bound]
\label{ass:uniform-perturbation}
For a fixed horizon $T$, the implemented MatrixPolicy and AdamW quotient iterates satisfy
\begin{align}
    q_{t+1}^{\MP,{\rm impl}}
    &=q_t^{\MP,{\rm impl}}-\eta_{\MP}P_{\MP}\nabla F(q_t^{\MP,{\rm impl}})+r_t^{\MP}, \\
    q_{t+1}^{\AW,{\rm impl}}
    &=q_t^{\AW,{\rm impl}}-\eta_{\AW}P_{\AW}\nabla F(q_t^{\AW,{\rm impl}})+r_t^{\AW},
\end{align}
with
\begin{equation}
    \norm{r_t^{\MP}}_2+\norm{r_t^{\AW}}_2\leq \delta
\end{equation}
for all $0\leq t<T$.  The residuals collect the effects of nonzero
$\epsilon_{\rm rms}$, AdamW denominator floors, EMA lag, clipping, decoupled weight decay,
and Taylor remainders of the nonquadratic objective.
\end{assumption}

\begin{proposition}[Stability of the strict gap]
\label{prop:complete-perturbation}
Fix $T\geq1$ and suppose the idealized dynamics satisfy a strict finite-time gap
\begin{equation}
    m_T
    :=F(q_T^{\AW})-F(q_T^{\MP})>0 .
    \label{eq:finite-gap}
\end{equation}
Assume the ideal trajectories remain in a compact region on which $F$ is Lipschitz and the
one-step maps are Lipschitz.  Then there exists $\delta_0>0$ such that, whenever
Assumption~\ref{ass:uniform-perturbation} holds with $\delta<\delta_0$,
\begin{equation}
    F(q_T^{\AW,{\rm impl}})-F(q_T^{\MP,{\rm impl}})>\frac{m_T}{2}>0 .
    \label{eq:stable-gap}
\end{equation}
\end{proposition}

\begin{proof}
Let $\Psi_{\MP}$ and $\Psi_{\AW}$ denote the ideal one-step maps.  By compactness, both
maps are Lipschitz on a neighborhood of the two ideal trajectories; let $C_\Psi$ be a
common Lipschitz constant.  A standard induction gives
\begin{equation}
    \norm{q_t^{\MP,{\rm impl}}-q_t^{\MP}}
    +
    \norm{q_t^{\AW,{\rm impl}}-q_t^{\AW}}
    \leq
    C_T\delta,
\end{equation}
where $C_T$ depends on $T$ and $C_\Psi$ but not on $\delta$.  Since $F$ is Lipschitz on the
same compact region, with constant $C_F$,
\begin{equation}
    \left|F(q_T^{\MP,{\rm impl}})-F(q_T^{\MP})\right|
    +
    \left|F(q_T^{\AW,{\rm impl}})-F(q_T^{\AW})\right|
    \leq
    C_F C_T\delta .
\end{equation}
Choose $\delta_0=m_T/(2C_FC_T)$.  Then $\delta<\delta_0$ implies
\begin{align}
    F(q_T^{\AW,{\rm impl}})-F(q_T^{\MP,{\rm impl}})
    &\geq F(q_T^{\AW})-F(q_T^{\MP})-C_FC_T\delta \\
    &>m_T/2.
\end{align}
This proves \eqref{eq:stable-gap}.
\end{proof}

In particular, if the ideal strict gap comes from
Corollary~\ref{cor:fixed-init-separation}, then the separation survives all sufficiently
small stabilizer, moment-lag, clipping, and weight-decay perturbations.  This is why the
exact homogeneous proof is still informative for the implemented optimizer: the conclusion
is finite-time stable as long as the run remains in the local regime described by the
assumptions.

\subsection{Statement of what is proved}
\label{app:complete-opt-proof:scope}

The result can be summarized as follows.  Under a local quotient model of the RLB block,
MatrixPolicy's rebalance step removes a nonfunctional positive scale degree of freedom and
thereby induces a canonical quotient preconditioner.  If that preconditioner has a strictly
smaller condition number than the quotient preconditioner induced by AdamW's raw-coordinate
adaptive state, then MatrixPolicy has a strictly smaller worst-case contraction factor on a
strongly convex quadratic quotient objective.  The finite-time gap is stable under small
implementation perturbations.

The theorem is a local quotient result for the mechanism tested empirically in
the main text: RLB exposes a gauge structure, and an optimizer that canonicalizes
and preconditions the quotient can converge faster than a coordinatewise optimizer
whose effective update depends on the arbitrary gauge representative.
